% 12-institutions.tex: the exchange fabric.
\section{Institutions: exchange, time and signatures without a message log}
\label{sec:institutions}

\subsection{The inversion}

A conventional exchange fabric between institutions makes an exchange provable by logging the
exchanged messages: a mediating node signs and timestamps the traffic so that a receipt can later be
shown to a third party. The evidence is real, and so is its cost: a store of who asked what about
whom, which is exactly the population-scale record the vocation forbids. The Polaris exchange fabric
is the inversion. It proves the same three things a message log proves, that an exchange occurred,
between these parties, and was authorized, while retaining none of the payload, by committing to the
request and the response by hash and never by content. Every primitive in it, signatures, the
detached verifier, the transparency logs, the trust graph, already existed; the fabric composes them.
\sImpl

\subsection{The registry: discovery, not configuration}

A consumer learns what an instance offers and trusts from one signed artifact: the protocol formats
and versions it speaks and its accepted algorithms, each service's kind, path template and
authentication, the transparency logs, every federated authority with its key register, the
contexts and the proof each requires, the in-context trust graph, and the relying parties by name
and scope only. Every fact is a view over the authority layer; the registry adds no truth of its own.
It must be signed by the active key it lists for its own publisher, so an impostor cannot publish one
in an authority's name, and it is non-transitive: it describes its publisher's instance, never
another's. The two-instance drill reads a service's path out of a fetched registry and calls it. No
token, no holder, no verification record appears: a map of the institutions, never of the people.
\sImpl

\subsection{The gateway and the receipt}

\figfile{gateway}{An institutional exchange through the gateway. The requester signs an envelope
under its registered key; the responder's own instance authenticates it by a key it already knows,
authorizes it through its own attestation in the envelope's context before anything leaves the
process, consumes the nonce in an append-only replay register, forwards the body only to an upstream
the operator configured, and returns the response with a signed receipt. Nothing but the receipt's
hash and the consumed nonce is written. The evidence of the exchange is the pair of envelope and
receipt, verifiable offline by a third party with no body.}{fig:gateway}

The receipt commits to the requester's key, the responder, the context, the two body hashes, which
attestation authorized the requester, and the instant. The mint endpoint accepts only hashes, so the
application cannot retain what it is never given. A third party verifies offline that the receipt is
authentic and that the requester was authorized in that context, given the manifests it trusts. A
party that holds a body can confirm the hash binds; a party that does not still holds the
responder's signed account. One correction from the outside review belongs here: a receipt alone
cannot prove the requester took part, because a dishonest responder can sign any receipt it likes;
the requester's own signed envelope beside it can, and the wire specification names the pair as the
evidence. The set of receipts is itself a transparency log, so an institution cannot later deny a
receipt existed and anyone can see how many it minted, while each receipt stays in the hands of its
parties. \sImpl

Authorization at the gateway is directional. Since \code{v9.333} the responding agency itself must
attest the requester's key in the envelope's context; an attestation by any other agency on the
instance authorizes nothing there. Before that fix the gateway accepted an attestation from anyone
on the instance, which was the transitive trust the federation layer refuses, reappearing one layer
up. \sImpl

\figfile{signing-chain}{A signed document and its long-term evidence. The document never leaves its
holder; its digest is signed by an institution, or by the holder's issuing authority on the
holder's behalf after a possession proof, recorded by credential hash. Evidence fixed at the instant
of signing, a timestamp from a distinct key, the signer's manifest, checkpoint and feed, together
with the verifier's own trust list, lets a verifier decide validity at that instant long after the
signing key has been retired. A signature whose instant falls after a compromise is refused even
though its embedded manifest says the key was active.}{fig:signing-chain}

\subsection{Time that learns nothing}

The timestamp authority binds an arbitrary SHA3-256 digest to an instant under an authority's
registered key. It accepts a digest, never content, and keeps no per-request record, because a
timestamp authority that logs every request is a store of who timestamped what and when. The
evidence is the signed statement in the requester's hands. Since \code{v9.341} a caller may choose
to anchor: the timestamp's own hash joins an append-only, witnessed log and the inclusion evidence
comes back stapled, so that a timestamp backdated under a stolen authority key is one absent from
every witnessed head of its claimed era. What the authority then retains is one digest and one
instant, a record of volume and timing, never of content, and the default request still leaves
nothing. \sImpl

\subsection{Signing a document so that it outlives the key}

A holder in Polaris holds a credential, not a signing key, so the model is notarial: the issuing
authority signs on behalf of a holder who has just proved possession, exactly as it asserts status
for one, and records who by hash. The authority's signature says that a holder of this credential,
active at this instant, asked it to sign this digest, and the embedded evidence lets anyone check
that claim later. \sImpl

Long-term validation is where the outside review and the repository's own drills sharpened the
claim. An authentic timestamp is not a trusted one, since anyone can sign one, and a signer's own
timestamp is backdatable by whoever holds the key. So the strong claim requires the timestamp to be
trusted by the verifier, from a key distinct from the signer's, active at the instant per the trust
list, and, if the verifier asks, anchored in a witnessed log or corroborated by a quorum of
independent authorities. A verifier given no timestamp anchors reports the facts and claims nothing.
\sImpl{} for the verifier; \sFut{} for anchor verification in the SDKs, which today only the detached
verifier performs.

\begin{layercard}{Layer card: institutions}
\qa{What it is}{The signed registry, the exchange gateway, the exchange receipt and its log, the timestamp authority with optional anchoring, document signing with long-term validation, the trust list.}
\qa{Why it exists}{Real software ecosystems need discovery, mediated requests, signatures and time. Each of those is conventionally built on a log of content, and the vocation forbids that log.}
\qa{What it trusts}{Registered institutional keys; the in-context trust graph; the verifier's own timestamp anchors and trust list.}
\qa{What it refuses}{A body at the mint endpoint; a URL named by a request; a replayed nonce; an attestation not made by the responder itself; a timestamp from the signer's own key as the strong evidence; a placeholder signature as authentication.}
\qa{What it can see}{Digests, keys, contexts, instants.}
\qa{What it cannot see}{The exchanged messages, the timestamped content, the signed document: none are ever given to it.}
\qa{If it fails}{A forged receipt fails the responder's signature; a dishonest responder's receipt is contradicted by the missing envelope; a stolen timestamp key is caught by the witnessed anchor or the quorum; a compromised signer key invalidates only what it signed after the instant.}
\qa{Checked by}{\code{check\_registry}, \code{check\_exchange\_gateway}, \code{check\_exchange\_receipt}, \code{check\_timestamp\_authority}, \code{check\_document\_signing}, \code{check\_ltv\_timestamp\_trust}, \code{check\_broker\_policy\_bound}.}
\qa{Now or later}{Now, between two instances. Real institutions are later.}
\qa{Inside and outside}{Inside: transparency, which watches the receipt and timestamp logs. Outside: the operations that keep all of it running.}
\end{layercard}

\nextlayer{Every layer so far is a claim about correctness. None of it matters if the system falls
over under load, loses a database, or cannot be deployed by anyone but its author. The next layer is
survival.}
